\documentclass{article}
\usepackage{graphicx} % Required for inserting images

\title{GCA}
\author{Alex Marinsek}
\date{February 2024}

\begin{document}

\maketitle

\section{Goals, learning outcomes, and expectations}

The goal of this course is to introduce students to the fundamentals of general purpose GPU programming. We shall dive into the concepts of parallel multi-thread processing, learning how to leverage the large number of GPU cores and how to efficiently use the device’s memory. Among others, concepts such as atomic operation, synchronous multi-threading, and dynamic parallelism are explored during the course. In addition to guided exercises, the students will partake in a short project, tackling a research or programming challenge of their liking using the newly acquired knowledge. By the end, the students are foreseen to have the necessary understanding of GPUs and parallel programming that will allow them to utilize GPUs in their future work. Throughout the course NVIDIA GPUs and C++ are employed (CUDA API).


\section{Lab sessions}

The below sections outline the six guided lab sessions (each 3h). This includes a brief summary, the goals/outcomes/expectations, and brief instructions-to-self.


\subsection{Getting acquainted}

An introduction to the GPU programming and execution workflow is provided. Emphasis is given on the distinction between host and device memory and functions. Moreover, the importance of identifying a reference point (index) within the parallel kernel instances is highlighted. They are given a simple task of flipping an array and debugging illegal memory operation using memcheck.


\subsubsection{Goals}
\begin{itemize}
    \item Know what is global and what is local scope
    \item Understand that memory has to be migrated between the two
    \item 
\end{itemize}

\subsubsection{Self-instructions}



\end{document}
